\documentclass[11pt]{article}
\usepackage[margin=1in]{geometry}
\usepackage[T1]{fontenc}
\usepackage{lmodern}
\usepackage{amsmath,amssymb,amsthm}
\usepackage{graphicx}
\usepackage[hidelinks]{hyperref}

\newtheorem{theorem}{Theorem}
\newtheorem{lemma}[theorem]{Lemma}
\newtheorem{remark}[theorem]{Remark}
\newcommand{\R}{\mathbb R}
\newcommand{\Lp}{L^p(\mu)}
\newcommand{\Var}{\operatorname{Var}_{\mu}}
\newcommand{\E}{\mathcal E}

\title{Contractive Centered $L^p$ Decay Forces Poincar\'e Without Symmetry}
\author{Run fa\_banach\_001, agent\_lane\_01}
\date{17 August 2026}

\begin{document}
\maketitle

\section*{Status and source question}

This is a \texttt{candidate\_full\_likely\_valid} result for Patrick
Cattiaux, Arnaud Guillin, and Cyril Roberto, \emph{Poincar\'e inequality and
the $L^p$ convergence of semi-groups}, arXiv:1003.0784 \cite{CGR}.
Section~4 on PDF page~12 discusses an invariant, not necessarily symmetric,
Markov diffusion semigroup.  The authors assume that, for some $p>2$, the
centered decay estimate
\begin{equation}\label{eq:Hp}
 N_p(P_t f)\le e^{-\lambda_p t}N_p(f),
 \qquad N_p(f)=\|f-\mu(f)\|_p,
\end{equation}
holds with the crucial prefactor $K_p=1$.  They obtain an $L^2$ exponential
estimate by duality and interpolation but state that, because they do not
know that $K_2=1$, they cannot conclude the Poincar\'e inequality.

\begin{center}
\includegraphics[width=0.98\linewidth]{figures/open_problem_crop.png}
\end{center}

The theorem below closes precisely this converse under the diffusion
hypotheses used by the source.

\section*{Setting and proof intuition}

Let $(P_t)_{t\ge0}$ be a Markov semigroup with invariant probability measure
$\mu$, generator $L$, and carr\'e du champ
\[
 \Gamma(f,g)=\frac12\bigl(L(fg)-fLg-gLf\bigr).
\]
Assume the source's diffusion chain rule, so
$\Gamma(\phi(f),\phi(f))=(\phi'(f))^2\Gamma(f,f)$ for smooth $\phi$.
Symmetry is not assumed.

Differentiating \eqref{eq:Hp} at zero gives the weighted inequality that the
source labels (2.4).  The tempting substitution
$h=\operatorname{sgn}(g)|g|^{2/p}$ makes its energy exactly quadratic, but
the source correctly notes two problems: this map is not $C^2$ and $h$ need
not have mean zero.

Both are removed at once.  Smooth the fractional power and first shift $g$
by a scalar $c_\varepsilon$ chosen so that the smoothed fractional power has
mean zero.  Its weighted energy multiplier is at most one, while its
$p$-th power converges to $|g-c|^2$.  The limiting $c$ need not be the mean:
variance is no larger than the squared distance from any constant.

\section*{Main result}

\begin{theorem}[nonsymmetric converse]\label{thm:main}
Suppose the invariant Markov diffusion semigroup above satisfies
\eqref{eq:Hp} for some $p>2$ and $\lambda_p>0$, for every $f\in L^p(\mu)$
and $t\ge0$.  Then
\begin{equation}\label{eq:poincare}
 \Var(g)\le \frac{p-1}{\lambda_p}
       \int\Gamma(g,g)\,d\mu
\end{equation}
for every $g$ in the diffusion core (and hence in the closed form domain).
In particular, centered $L^2$ convergence holds with prefactor one and rate
at least $\lambda_p/(p-1)$.
\end{theorem}

\begin{proof}
It is enough to work with real-valued core functions.  First let $h$ be a
mean-zero core function.  Raising \eqref{eq:Hp} to the $p$-th power and
taking the right derivative at $t=0$ gives
\[
 p\int |h|^{p-2}hLh\,d\mu
 \le -p\lambda_p\int|h|^p\,d\mu.
\]
Invariance and the diffusion chain rule applied to $|h|^p$ give
\[
 -\int |h|^{p-2}hLh\,d\mu
 =(p-1)\int|h|^{p-2}\Gamma(h,h)\,d\mu.
\]
Thus
\begin{equation}\label{eq:weighted}
 \lambda_p\int|h|^p\,d\mu
 \le(p-1)\int|h|^{p-2}\Gamma(h,h)\,d\mu.
\end{equation}

Fix a bounded real core function $g$; the general form-domain statement
follows by the standard truncation and closure argument.  Put
$\alpha=2/p\in(0,1)$ and, for $\varepsilon>0$, define the smooth odd function
\[
 \psi_\varepsilon(s)
 =s(s^2+\varepsilon)^{(\alpha-1)/2}.
\]
It is strictly increasing because
\begin{equation}\label{eq:derivative}
 \psi_\varepsilon'(s)
 =(s^2+\varepsilon)^{(\alpha-3)/2}
   (\varepsilon+\alpha s^2)>0.
\end{equation}
The continuous decreasing function
\[
 c\longmapsto\int\psi_\varepsilon(g-c)\,d\mu
\]
has opposite signs at the essential lower and upper bounds of $g$.
Choose a zero $c_\varepsilon$; these centers stay in that bounded interval.
Then
$h_\varepsilon=\psi_\varepsilon(g-c_\varepsilon)$ has mean zero and belongs
to the diffusion core by smooth functional calculus.

Two elementary scalar estimates contain the whole argument.  Since
$\alpha p=2$ and $\alpha-1<0$,
\begin{equation}\label{eq:Lpbound}
 |\psi_\varepsilon(s)|^p
 =|s|^p(s^2+\varepsilon)^{(\alpha-1)p/2}
 \le |s|^2.
\end{equation}
Also, \eqref{eq:derivative} implies
$0<\psi_\varepsilon'(s)\le
(s^2+\varepsilon)^{(\alpha-1)/2}$, whence
\begin{align}
 |\psi_\varepsilon(s)|^{p-2}
 |\psi_\varepsilon'(s)|^2
 &\le |s|^{p-2}
       (s^2+\varepsilon)^{(\alpha-1)p/2} \notag\\
 &=\left(\frac{s^2}{s^2+\varepsilon}\right)^{(p-2)/2}
 \le1.\label{eq:energybound}
\end{align}

Apply \eqref{eq:weighted} to $h_\varepsilon$.  The diffusion chain rule and
\eqref{eq:energybound} yield
\begin{equation}\label{eq:beforelimit}
 \lambda_p\int|\psi_\varepsilon(g-c_\varepsilon)|^p\,d\mu
 \le(p-1)\int\Gamma(g,g)\,d\mu.
\end{equation}
Choose a sequence $\varepsilon\downarrow0$ along which
$c_\varepsilon\to c$.  Pointwise,
\[
 |\psi_\varepsilon(g-c_\varepsilon)|^p
 \longrightarrow |g-c|^2.
\]
The bound \eqref{eq:Lpbound}, together with boundedness of $g$ and the
centers, permits dominated convergence in \eqref{eq:beforelimit}.  Therefore
\[
 \lambda_p\int|g-c|^2\,d\mu
 \le(p-1)\int\Gamma(g,g)\,d\mu.
\]
Finally,
$\Var(g)=\inf_{a\in\R}\int|g-a|^2d\mu\le\int|g-c|^2d\mu$, proving
\eqref{eq:poincare}.

For completeness, invariance also gives
\[
 \frac{d}{dt}\Var(P_tg)
 =-2\int\Gamma(P_tg,P_tg)\,d\mu.
\]
Combining this identity with \eqref{eq:poincare} and Gronwall's lemma gives
\[
 \Var(P_tg)\le
 e^{-2\lambda_p t/(p-1)}\Var(g),
\]
which is the asserted $K_2=1$ conclusion.
\end{proof}

\section*{Verification, scope, and novelty status}

The included scalar script checks \eqref{eq:derivative},
\eqref{eq:Lpbound}, and \eqref{eq:energybound} on 1,159,968 logarithmically
distributed triples $(p,\varepsilon,s)$; all cases pass.  This is a regression
check, not a substitute for the displayed proof.

The hypothesis $K_p=1$ is essential to this route: if a prefactor greater
than one is present, differentiation at zero does not yield
\eqref{eq:weighted}.  The proof is specific to diffusions through the exact
chain rule and does not assert a pure-jump analogue.  It requires no symmetry,
duality, or interpolation.

Eight focused attempts are recorded in the accompanying attempt note.  The
first interpolation and finite-state operator-norm routes retained an
uncontrolled $K_2$.  Direct differentiation isolated \eqref{eq:weighted}; the
naive fractional power reproduced the source's centering obstruction.  The
successful upgrade was the shifted smooth fractional power above.  Further
checks established the energy bound, closure step, quantitative constant,
and nonsymmetric $L^2$ consequence.

Cheap run indexes and bounded web/arXiv searches on 17 August 2026 used the
exact source sentence and the phrases \emph{$K_p=1$ Poincar\'e nonsymmetric
semigroup}, \emph{centered $L^p$ exponential convergence Poincar\'e}, and the
weighted formula \eqref{eq:weighted}.  They found the arXiv and published
versions of \cite{CGR} and later related centered Sobolev results, but no
resolution of this exact converse.  Novelty confidence is moderate-high
pending expert review.

\paragraph{Human-review recommendation.}
Check the right-differentiation/core step leading to \eqref{eq:weighted} and
the routine truncation/closure extension.  The scalar regularization and
limiting argument are otherwise self-contained.

\begin{thebibliography}{9}
\bibitem{CGR}
P.~Cattiaux, A.~Guillin, and C.~Roberto,
\emph{Poincar\'e inequality and the $L^p$ convergence of semi-groups},
Electron. Commun. Probab. 15 (2010), 270--280; arXiv:1003.0784.
\end{thebibliography}

\end{document}
